\documentclass[11pt, a4paper]{article}

\usepackage{fontspec}
\usepackage{geometry}
\usepackage{fancyhdr}
\usepackage[hidelinks]{hyperref}
\usepackage[normalem]{ulem}
\usepackage{multicol}
\usepackage{enumitem}

\geometry{
  top=2cm,
  bottom=2cm,
  left=2cm,
  right=2cm,
  marginparsep=4pt,
  marginparwidth=1cm
}

\renewcommand{\headrulewidth}{0pt}
\pagestyle{fancyplain}
\fancyhf{}
\lfoot{}
\rfoot{}

\setlength{\parindent}{0pt}
\setlength{\parskip}{0pt}

\usepackage{xunicode}
\defaultfontfeatures{Mapping=tex-text}

\setlist[itemize]{leftmargin=1.4em, itemsep=0pt, topsep=1pt, parsep=0pt}

\newcommand{\code}[1]{\texttt{#1}}

\begin{document}

\small

\begin{center}
\normalsize\textbf{DSST389: Exam I --- Reference Sheet}
\end{center}

\smallskip

\textit{This sheet is attached to your exam. You do not need to memorize any of it.}

\smallskip

\textbf{Setup.}
Importing \code{funs} registers a namespace called \code{geo} on every Polars DataFrame.
Geometries are stored as Well-Known Text (WKT) in a column named \code{geometry}; the
coordinate system is recorded in a column named \code{\_crs}.

\smallskip

\textbf{Reading and building geometries.}

\begin{itemize}
  \item \code{funs.read\_file(path)} --- read a GeoJSON, Shapefile, or GeoPackage into a
        Polars DataFrame with the geometries converted to WKT.
  \item \code{df.geo.points\_from\_xy(x="lon", y="lat", crs=4326)} --- build a point
        \code{geometry} column from two coordinate columns and record the CRS.
  \item \code{df.geo.to\_crs(code)} --- reproject the geometries into the given EPSG code.
  \item \code{df.geo.explore(column=None, tooltip=["name"])} --- interactive map.
\end{itemize}

\smallskip

\textbf{Metrics.}
Some of these are properties and take no parentheses. A property that returns
\textbf{geometries} overwrites the \code{geometry} column; one that returns \textbf{numbers}
appends a new column named after itself.

\begin{itemize}
  \item \code{df.geo.centroid} --- \textit{property.} Replaces each geometry with its
        center point.
  \item \code{df.geo.area} --- \textit{property.} Appends a column \code{area}, in squared
        CRS units.
  \item \code{df.geo.x} \quad \code{df.geo.y} --- \textit{properties.} Append a column
        \code{x} or \code{y} holding a point's coordinate.
  \item \code{df.geo.buffer(distance)} --- replaces each geometry with the region within
        \code{distance} of it, measured in CRS units. A point becomes a disc.
  \item \code{df.geo.count\_coordinates()} --- appends a column \code{count\_coordinates}
        giving the number of vertices in each geometry.
\end{itemize}

\smallskip

\textbf{Spatial joins.}
Both tables must be in the same CRS. The output keeps the geometry of the \textbf{left}
table. Columns whose names appear in both tables get \code{\_left} and \code{\_right}
suffixes, and an \code{index\_right} column is added. The join is an \textbf{inner} join:
left rows with no match are dropped, without an error.

\begin{itemize}
  \item \code{df.geo.sjoin(other, predicate="intersects")}
  \item \code{df.geo.sjoin\_nearest(other, exclusive=False, distance\_col=None)} --- matches
        each left row to the closest right row. \code{exclusive=True} forbids a row from
        matching itself; \code{distance\_col} names a column to hold the distance, in CRS
        units.
  \item Predicates: \code{intersects} (share any space, including a boundary),
        \code{contains} (the first fully encloses the second), \code{within} (the first is
        fully enclosed by the second), \code{touches} (share a boundary but no interior).
\end{itemize}

\smallskip

\textbf{Changing the geometries.}

\begin{itemize}
  \item \code{df.geo.dissolve(by="col", aggfunc=\{"area":"sum"\})} --- the spatial
        counterpart of \code{group\_by} and \code{agg}. Merges the geometries within each
        group into one and erases the internal boundaries. The group key comes back as an
        ordinary column.
  \item \code{df.geo.explode()} --- one row per part of a multi-part geometry. The other
        columns are copied onto every part.
  \item \code{df.geo.simplify(tolerance)} --- drops vertices, keeping the boundary within
        \code{tolerance} of its original position, in CRS units. Each geometry is
        simplified independently, so shared borders can develop gaps: simplify for display,
        not for analysis.
\end{itemize}

\smallskip

\textbf{Plotting.}
\code{geom\_map} draws the \code{geometry} column and takes \code{fill}, \code{color}, and
\code{alpha} as either fixed values or aesthetics. \code{coord\_fixed} keeps one unit on the
x-axis the same length as one unit on the y-axis.

\begin{verbatim}
(
    ggplot(state, aes(fill=c.area))
    .geom_map(color="white")
    .coord_fixed()
)
\end{verbatim}

\smallskip

\textbf{Coordinate reference systems.}
Areas and distances come back in the units of the current CRS. In a meter-based CRS an area
is in square meters, and $1 \textrm{ km}^2 = 1{,}000{,}000 \textrm{ m}^2$.

\smallskip

\sloppy
\begin{itemize}
  \item \textbf{EPSG:4326}, WGS 84 --- longitude and latitude, in \textbf{degrees}. GPS,
        raw data, and the default of \code{points\_from\_xy}. Never for areas or distances.
  \item \textbf{EPSG:5069}, NAD83 Conus Albers --- equal-area, contiguous United States.
        Units: \textbf{meters}.
  \item \textbf{EPSG:3857}, Web Mercator --- screen maps and Google Maps. Badly distorts
        area.
  \item \textbf{EPSG:32618}, UTM Zone 18N --- the US East Coast. Units: meters.
\end{itemize}
\fussy

\end{document}
